\section{Validation Protocols}
\subsection{Literature-Calibrated Physical Model}
The vehicle is a traceable cross-source heavy-duty composite, not an experimentally identified individual truck. Mass, braking force, thermal capacity, and cooling follow the National Highway Traffic Safety Administration (NHTSA) heavy-truck S-cam report \cite{nhtsa2010scam}; total length follows the American Association of State Highway and Transportation Officials (AASHTO) design vehicle \cite{aashto2004}; resistance and drag use a Class-8 study \cite{wang2015fuel}; acceleration and friction-lag bounds use heavy-truck records \cite{vegamoor2018,yanakiev1998}; the auxiliary lag and digitized retarder map use published heavy-truck models \cite{li2026retarder,donkers2020powertrain}; the representative fade curve is digitized from a National Transportation Safety Board (NTSB) study \cite{ntsb2002simulation}; ambient and the independent thermal-discrepancy model use \cite{yan2018thermal}; and communication ranges use field measurements \cite{hu2021delay,gao2016dsrc}. The 10\% and 1-km descent benchmark follows \cite{devika2022fade}. Table~\ref{tab:final-parameters} is generated from the provenance registry.

\begin{table*}[!t]
\centering\scriptsize
\caption{Selected literature-calibrated parameters; full intervals and source-level provenance are archived with the configuration.}
\label{tab:final-parameters}
\begin{tabular}{lllll}
\toprule
Parameter & Symbol & Nominal & SI unit & Provenance \\
\midrule
Vehicle mass & $m$ & 33200 & kg & DIRECT \\
Maximum friction force & $\bar u^f$ & 123603 & N & DERIVED \\
Friction actuator lag & $\tau_f$ & 0.14 & s & DIRECT \\
Auxiliary actuator lag & $\tau_a$ & 0.25 & s & DERIVED \\
Lumped thermal capacity & $C$ & 15439.7 & J/K & DERIVED \\
Cooling coefficient & $k$ & 624.334 & W/K & DERIVED \\
Critical brake temperature & $T_{\rm crit}$ & 463.708 & K & DERIVED \\
Ambient temperature & $T_a$ & 298.15 & K & DIRECT \\
Nominal residual heat & $q_w$ & 0 & W & DERIVED \\
\bottomrule
\end{tabular}
\end{table*}


The validation campaign contains ten studies specified before evaluation: M1 physical feasibility, M2 long-descent allocation, M3 hot-brake conflict, M4 predictive warning, M5 horizon dependence, M6 endpoint and joint design sensitivity, M7 delayed communication, M8 low-speed branch behavior, M9 numerical string stability, and M10 computation time. All 92 campaign runs passed integrity checks. The 20 joint M6 samples are a seeded design-sensitivity set, not draws from a physical probability distribution. A separate follow-up contains 81 boundary-approach scenarios plus the M7 diagnosis; no physical parameter or supervisor threshold was retuned.

For the boundary study, we evaluate the same instantaneous hard rows on a $25\times51$ grid of initial temperature (365--461~K, 4-K increments) and bumper gap (15--90~m, 1.5-m increments) at 15~m/s and a 10\% downgrade. Friction-only feasibility fixes $u_i^a=0$ in the same two-command geometry; dual-brake feasibility retains its physical auxiliary interval. This is an initial command-domain comparison, not a closed-loop campaign. A same-plant supervisor ablation evaluates three 10-s, three-truck scenarios with identical nominal allocator, physical limits, initial states, leader trace, channel, and sampling; the comparator disables only predictive/upstream triggering while retaining instantaneous hard-QP projection. Fixed 50-, 100-, 200-, and 300-ms channel cases use the M7 plant and policy to evaluate head-truck reserve deviation and critical attribution against an instantaneous centralized oracle. The 300-point stratified design varies only nondegenerate declared parameter intervals and the fade digitization offset around one boundary state; it is a deterministic sensitivity design, not a probability law. Source CSV rows and generating scripts accompany these analyses.

The predictive evidence is organized by: (P1) instantaneous versus horizon certificate, (P2) vehicle--time certificate map, (P3) centralized versus recursively available fleet bottleneck, and (P4) physical-limit attribution. Only evaluated, provenance-checked results are reported; planned but unexecuted plots are not treated as evidence.

\subsection{Secondary Matched Learning Protocol}
The ten-seed DiffQP/StopGradient diagnostic matches initialization, random streams, scenario order, reward, optimization budget, forward QP, executed action, and evaluation seeds. Each run uses 20 updates, 12 steps per update, two policy-optimization epochs, and three trucks. Its only intended difference is whether the provisional QP output contributes a valid implicit derivative to the intervention objective; nonregular samples use a zero-QP-gradient fallback. The full protocol and paired outcomes remain in the reproducibility package, because learning performance is not the feasibility contribution.

\subsection{External Recent T-ITS Baseline}
We independently implemented the mixed-autonomy safe-MARL/cooperative-CBF method of Zhou \emph{et al.} \cite{zhou2026cooperative}. The native reimplementation reproduced collision-free behavior and nonnegative cooperative-CBF values in the three source-style disturbance families after 450 episodes, 450,000 environment steps, and 439 PPO updates. Average connected-and-automated-vehicle (CAV) headway was 1.385, 1.149, and 1.144~s; Average Absolute Velocity Error (AAVE), using the source-defined metric name, was 1.620, 1.362, and 4.699~m/s; minimum CBF values were $6.08\times10^{-10}$, $8.32\times10^{-10}$, and $1.00\times10^{-10}$; and every collision count was zero. These results establish qualitative and functional agreement of the independent implementation, not numerical reproduction of every published value. Architecture and optimization details not reported by the source were specified before formal evaluation.

For the matched transfer, the same trained baseline produces a scalar safe acceleration on the unchanged heavy-duty longitudinal/thermal plant. A fixed auxiliary-first adapter converts the implied total braking demand to $(u_i^f,u_i^a)$ by using available auxiliary braking first, assigning the remainder to friction braking, and enforcing only basic magnitude, gear, and power limits. It does not use $\rho_i$, $\rho_i^{H,\mathrm{cert}}$, $\rho_i^{H,\mathrm{up}}$, the thermal HOCBF, fade CBF, predictive tube, certificate supervisor, or physical-limit attribution. Certificate quantities reported for this baseline are post-hoc common feasibility evaluators and never feed back to its control.

Six deterministic cases were specified before evaluation: nominal long descent; three hot starts; emergency braking during descent; and nominal communication delay. Each pair shares the plant, initial condition, road, leader, ambient condition, communication trace, integration settings, duration, and seed 20260918. No per-case retuning, stochastic replication, or confidence interval is used. The heavy-duty experiment is an adapted Zhou \emph{et al.} baseline, not the original physical environment of Zhou \emph{et al.}

\subsection{Implementation and Integrity Checks}
The detailed mapping is archived in the theory-to-validation traceability matrix. Collision, thermal, fade, and low-speed rows in (11)--(24) are checked by symbolic differentiation, affine/QP agreement, and hot/cold tests at $v=0$, $v\approx0$, $v=v_\epsilon$, and $v>v_\epsilon$. The directional margin in (25) is checked by interval-endpoint regression, not a formal rounding proof; the hard QP and exact certificate in (26)--(33) by forward matching, hard-row residuals, and independent linear-program (LP) vertex comparisons; and the predictive construction in (34)--(37) by the $H=1$ reduction, sampled interval enclosure checks on 5,000 states, horizon/uncertainty monotonicity, and targeted boundary-approach cases. Fleet/recursive (38)--(40), learning/KKT (41)--(42), and supervisor (43) map to min/argmin/delay, residual/gate/finite-difference/fallback, and changed-action/final-row tests, respectively. Low-speed intersection (44), sampled-data rate bound (45), and low-speed ZOH error bound (46) map to dedicated branch/intersection, executable rate-bound, and executable error-bound tests, respectively.

These checks detect algebraic and coding regressions and test consistency of the executable formulas under declared bounds. They do not identify real uncertainty bounds by themselves, prove arbitrary continuous-time containment or global backup-domain nonemptiness, or replace real-vehicle, field, or hardware-in-the-loop validation.
